\subsection{Complete full and metabelian sections over two families of arithmetic sources}

\begin{statement}{E.104}{Proved}{All cubic sections over the two-prime source families of E.55 and E.56}
Keep the actual based cubic geometric pro-three group $N$,
its intrinsic metabelian quotient $M=N/N''$, the original
complex conjugation $c$, and its genuine action $J$.
Define the following explicit sets of rational primes:
\[
 \begin{aligned}
 \mathscr I&=\{\ell\text{ prime}:\ell\equiv2,5\pmod9\},\\
 \mathscr P&=\{\ell\text{ prime}:v_3(\ell-1)=1,
                      \ 3^{(\ell-1)/3}\not\equiv1\pmod\ell\}.
 \end{aligned}
 \tag{E104.1}
\]
For every $\ell\in\mathscr I\cup\mathscr P$, let
$G_{3\ell}$ be the actual relative pro-three arithmetic
group over $K=\mathbb Q(\zeta_3)$, unramified outside
the primes over three and $\ell$, with its retained
order-two quotient. Write
\[
 \mathscr S_{3\ell}(H)=H^1_{\rm cts}(G_{3\ell},H)
       \quad(H=N,M)
 \tag{E104.2}
\]
for geometric-conjugacy classes under the genuine based
action. Inflation along the specified arithmetic quotient
identifies this with the corresponding locus of global
sections. The following coordinates compute the entire
space (E104.2), and every full-to-metabelian fibre.

For $\ell\in\mathscr I$, put $q=\ell^2$.
There are actual marked generators
\[
 \begin{gathered}
 G_{3\ell}=(D_q\mathop{*}_3\langle\rho\rangle)
                       \rtimes\langle c\rangle,
 \quad D_q=\langle t,\sigma\mid\sigma t\sigma^{-1}=t^q\rangle_3,\\
 ctc=t^{-1},\quad c\sigma c=\sigma,\quad c\rho c=\rho^{-1}.
 \end{gathered}
 \tag{E104.3}
\]
In the original Kummer marking their columns are
\[
 \begin{array}{c|ccc}
  &t&\sigma&\rho\\\hline
  \chi_3&1&q&1\\
  \kappa_3(\ell)&1&0&0\\
  \kappa_3(3)&0&0&2
 \end{array}.
 \tag{E104.4}
\]
Let $T=\Phi_\tau$ be the original arithmetic operator
of E.37, and set $F_\ell=\Phi_\sigma$.
The free generator is chosen so that
\[
 \Phi_t=1,\quad \Phi_\rho=TJT^{-1}J,
 \quad\mathcal L:=\Phi_\rho J=TJT^{-1}.
 \tag{E104.5}
\]
Use the same three rational cusps, rational tangents
and $J$-fixed paths $p_i$ as in E.100, and put
$u_{i,\ell}=p_iF_\ell(p_i)^{-1}$.
Every ramified section has a unique representative
\[
 \begin{array}{c|cccc}
    &t&\sigma&\rho&c\\\hline
  a_{i,r,y}&d_i^r&u_{i,\ell}&T(y)&1
 \end{array},\qquad
 i\in\{O,A_0,B_0\},\quad
 r\in\mathbb Z_3\setminus\{0\},\quad J(y)=y^{-1}.
 \tag{E104.6}
\]
Uniqueness here means uniqueness after the unique
$J$-fixed geometric conjugator aligning its tame
restriction with $(d_i^r,u_{i,\ell})$.

For $\ell\in\mathscr P$, the marked source is instead
\[
 \begin{gathered}
 G_{3\ell}=(D_\ell\mathop{*}_3D'_\ell)
                          \rtimes\langle c\rangle,\\
 D_\ell=\langle t,\sigma\mid\sigma t\sigma^{-1}=t^\ell\rangle_3,
 \quad c\sigma c=\sigma',\quad ctc=(t')^{-1}.
 \end{gathered}
 \tag{E104.7}
\]
Here the exact root-normalized generators are those
of E56.11--13. Put $T_\ell=\Phi_\sigma$.
For each of the nine $K$-rational cusps $D$, use its
actual positive peripheral element $\delta_D$ and
path value $z_{D,\ell}$, so that
\[
 z_{D,\ell}T_\ell(\delta_D)z_{D,\ell}^{-1}
                         =\delta_D^\ell.
 \tag{E104.8}
\]
The complete ramified coordinates are
\[
 (D,r,b,[g])\in
 \{\text{nine cusps}\}\times
 (\mathbb Z_3\setminus\{0\})\times\{0,1,2\}
                         \times H^J\backslash H,
 \tag{E104.9}
\]
with actual values
\[
 \begin{gathered}
 a(t)=g\delta_D^r g^{-1},\qquad
 a(\sigma)=g\delta_D^b z_{D,\ell}T_\ell(g)^{-1},\qquad a(c)=1,\\
 a(t')=J(a(t))^{-1},\qquad a(\sigma')=J(a(\sigma)).
 \end{gathered}
 \tag{E104.10}
\]
Every row in (E104.6) and (E104.10) is a global section
in its stated source; every geometric stabilizer is trivial.

In both families the zero-inertia part is exactly
$\mathscr S_3(H)=H^1_{\rm cts}(\Gamma,H)$ of E.41.
The actual full-to-metabelian map is onto on the whole
space. It preserves the displayed cusp, exponent and
finite residue coordinates. Its ramified fibres are
\[
 \prod_{m\ge4}\mathbb Z_3^{\nu_m(N)-\nu_m(M)},
 \tag{E104.11}
\]
with the exact all-degree ranks and mutually inverse
Hall maps of E100.7 and E98.30--31. Every zero-inertia
fibre is
\[
 \prod_{m\ge4}\mathcal C_m(N''),\qquad
 \mathcal C_m(N'')=K_m/(T_m-I)K_m^J,
 \tag{E104.12}
\]
with the actual nonlinear comparison maps E100.23--26.
In particular its first factor is exactly
$\mathbb Z_3^{504}\oplus C_3^{144}\oplus C_9^{56}
\oplus C_{27}^{30}\oplus C_{81}^{9}$ by E.102.
No full ramified class maps to a zero-inertia target.
The whole space retains the quotient topology and
the $r\to0$ boundaries described in the proof.

The arithmetic source conditions in (E104.1) specify
the scope of this classification. E.23 continues to
retain the unrestricted global source and the other
requested geometric prime components. The manuscript
anchor is \S51.16, folios733v--736v, scans172--178,
formulas(371)--(379), Part E(c),(d) and follow-up4.
\end{statement}

\begin{proof}
\emph{Original-marking transport for every inert prime.}
E.55 constructs the actual tame factor at the unique
$K$-prime over $\ell$, with residue cardinality $q$.
We specify its transport to the original $c$ and roots.
Let its local order-two element be $d=g c g^{-1}$.
For either original rational Kummer cocycle $\kappa$
we have $\kappa(c)=0$ and
$\kappa(d)=2\kappa(g)$. Multiplying the root system
by the compatible roots of unity with exponent
$e=\kappa(g)$ changes the cocycle to
$\kappa+(\chi_3-1)e$ and makes those roots $d$-fixed.
The local Frobenius $\widehat\sigma$, chosen as
in E.55, fixes these roots. Therefore
\[
 \begin{gathered}
 \kappa(\widehat\sigma)=-(q-1)\kappa(g),\\
 \kappa(g^{-1}\widehat\sigma g)
 =\chi_3(g)^{-1}
     \bigl(\kappa(\widehat\sigma)+(q-1)\kappa(g)\bigr)=0.
 \end{gathered}
 \tag{E104.13}
\]
Conjugate the entire source by $g^{-1}$ and raise
its tame inertia to the unit power $\chi_3(g)$.
These operations have the stated inverse conjugation
and inverse unit power, preserve the tame relation,
and give the first two columns of (E104.4) for the
original roots and the original $c$.

Lift the original $\tau$ to this arithmetic source.
Multiply its lift on the left by the negative of its
$\kappa_3(\ell)$-power of $t$, obtaining $\gamma$.
Its image in the unramified quotient is still $\tau$.
The actual based geometric action is unramified at
$\ell$ by E.34 and E.78, so $\Phi_\gamma=T$.
Set $\rho=\gamma c\gamma^{-1}c$. The cocycle identities
give its last column in (E104.4), the exact relation
$c\rho c=\rho^{-1}$ and (E104.5).
The two Frattini characters of $\sigma,\rho$ in the
free rank-two unramified group $P$ have matrix
\[
 \begin{pmatrix}(q-1)/3&0\\0&2\end{pmatrix}\pmod3.
 \tag{E104.14}
\]
Its determinant is a unit because $v_3(q-1)=1$.
Thus these are an actual free basis of $P$.
The inverse basis map is the compatible inverse on
finite characteristic quotients, by the finite-index
Hopfian argument in E.100.

The decomposition factor and this particular $\rho$
give (E104.3), including both inverse maps.
Indeed the kernel of its map to $P$ is the continuous
free product of the actual tame inertia groups over
$P/\langle\sigma\rangle$: on the arithmetic side this
is the opened Wingberg theorem
\cite[Lemma1.1, Theorem2.4 and their proofs]{E_WingbergRiemann};
on the displayed source it is its pro-three coset tree.
The generator inclusions give the identity on $P$
and the same isomorphism on every inertia factor.
The universal properties give inverse maps on the
kernel, hence on the whole source. Choosing the same
generators in the original arithmetic group gives
the reverse continuous section of its projection.
Their conjugation relations make this map $c$-equivariant.
This proves the required source maps for every
$\ell\in\mathscr I$, with all original markings retained.

\emph{The actual split-prime sources.}
For $\ell\in\mathscr P$, E.56 constructs actual
Frobenii with cyclotomic character $\ell$ and Kummer
values of three equal to $1,-1$. Its explicit
conjugator E56.12 changes the original cubic value
$t_0\equiv1\pmod3$ to exactly one; it preserves
the original complex conjugation and root systems.
The two unramified Frattini columns have determinant
$-2(\ell-1)/3$ modulo three, hence give a free basis
of the same $P$. The two actual decomposition groups
then have the free-product presentation (E104.7).
For clarity, its kernel comparison uses the two
induced inertia bundles with relations
$(p\sigma^z,t^u)\sim(p,t^{\ell^zu})$ and the conjugate
relations. The coset-tree inclusions and the actual
arithmetic inertia inclusions are inverse on those
bundles and on $P$. This is the precise E.56
source isomorphism and its reciprocal arithmetic section.

For both families, the genuine action on $N$ factors
through the original unramified $\Gamma=P\rtimes\langle c\rangle$.
Over $K$ its image is pro-three. A cocycle paired
with this action is therefore a homomorphism to a
pro-three semidirect product; it factors through
the displayed relative pro-three source. Inflation
along the actual arithmetic quotient is the inverse
construction. A geometric conjugacy after inflation
already conjugates the values on its onto quotient,
so inflation is injective on the stated classes.

\emph{The exact section equations and their topology.}
The order-two equation for a cocycle value $v$ at
$c$ is $vJ(v)=1$. Conjugation by $v^{-1/2}$ makes
that value one; all remaining conjugators are $H^J$.
For the inert family, write its other values as
$(\nu,u,x)$ at $(t,\sigma,\rho)$. The complete
normalized space is the actual quotient
\[
 \left\{(\nu,u,x):
 J(\nu)=\nu^{-1},\ J(u)=u,\
 \mathcal L(x)=x^{-1},\ uF_\ell(\nu)u^{-1}=\nu^q
 \right\}/H^J .
 \tag{E104.15}
\]
The action is
$(\nu,u,x)\mapsto
(h\nu h^{-1},huF_\ell(h)^{-1},hx\Phi_\rho(h)^{-1})$.
For the split family the complete normalized space is
\[
 \{(\nu,u):uT_\ell(\nu)u^{-1}=\nu^\ell\}/H^J,
 \quad h(\nu,u)=(h\nu h^{-1},huT_\ell(h)^{-1}).
 \tag{E104.16}
\]
The primed values are forced by $c$, and the second
tame equation is exactly the $J$-image of the first.
The free-product universal property proves sufficiency
as well as necessity of these equations. Their solution
sets are closed in compact powers of $H$; the compact
geometric actions give Hausdorff quotient spaces.
These are the topologies used throughout the statement.

\emph{Every inert ramified class.}
E.61 and E.66 classify the complete ramified full
and metabelian local sections at every good prime
away from three. Their maps apply to the actual
local factor $D_q$, with all nonunit exponents
included. The nine local labels are
$(D,r,\bar b)$, with $r\ne0$ and
$\bar b\in\mathbb Z_3/(q-1)\mathbb Z_3=\mathbb F_3$.
Choose conjugate paths for each nonrational cusp
pair, and the $J$-fixed paths of E.100 at the three
rational cusps. The actual involution on these
packets is $(D,r,\bar b)\mapsto(\bar D,r,-\bar b)$:
its source inertia inversion cancels its peripheral
inversion, whereas $\sigma$ is fixed and the
peripheral factor in its value is inverted.
The fixed labels are exactly $O,A_0,B_0$ and
$\bar b=0$. This derives the descent labels
from the actual source equations.

The rational path torsors factor through $\Gamma$
by the centre-free endpoint argument of E.100.
Thus their $\sigma$-values are precisely
$u_{i,\ell}$, with
$J(u_{i,\ell})=u_{i,\ell}$ and
$u_{i,\ell}F_\ell(d_i)u_{i,\ell}^{-1}=d_i^q$.
The local class has an aligning conjugator to
$(d_i^r,u_{i,\ell})$. Its ambiguity centralizes
$d_i^r$, hence equals $d_i^k$, by the exact full
and metabelian centralizer maps of E.61 and E.66.
Preserving the Frobenius value forces $k=qk$,
so $k=0$. Uniqueness also makes the conjugator
$J$-fixed, since its $J$-image aligns the same
two normalized sections. The remaining free value
is exactly $x\in\operatorname{Anti}_{\mathcal L}(H)$.
The literal conjugacy $\mathcal L=TJT^{-1}$ gives
$x=T(y)$ with $y\in\operatorname{Anti}_J(H)$.
This proves (E104.6) and its inverse, including
triviality of the full geometric stabilizer.

\emph{Every split ramified class.}
Apply the same complete local classification to
$D_\ell$. For each cusp its whole path torsor
is unramified outside three: its endpoint is obtained
from the base tangent by a coordinate permutation
and multiplication by third roots of unity over $K$.
The centre-free endpoint argument therefore gives
the actual value $z_{D,\ell}$ in (E104.8).
The local inertia and Frobenius values are exactly
the first line of (E104.10), initially with
$b\in\mathbb Z_3$. For fixed $D,r$, their entire
right ambiguity is
\[
        (g,b)\longmapsto(g\delta_D^k,b+(\ell-1)k).
 \tag{E104.17}
\]
Substitution using (E104.8) proves this equality
of both actual cocycle values. As
$(\ell-1)\mathbb Z_3=3\mathbb Z_3$, there is a
unique such $k$ making $b$ its representative
in $\{0,1,2\}$. The residual global conjugations
are exactly left multiplication by $H^J$.
The inertia abelianization determines $D,r$;
the centralizer proof and (E104.17) determine
the rest uniquely. This proves the maps
(E104.9)--(E104.10). A stabilizer would have
the same exponent equation $k=\ell k$, so is trivial.

The quotient $H^J\backslash H$ has the actual
inverse coordinates
$[g]\mapsto g^{-1}J(g)$ and
$x\mapsto[x^{-1/2}]$ from E98.10.
Its Hall coordinates and those of the inert
anti set are the integral ones of E.100.
On each shell $v_3(r)=e$ the parameter maps
are continuous bijections from compact spaces
to Hausdorff spaces and hence homeomorphisms.
For fixed remaining parameters, $r\to0$ in
(E104.6) or (E104.10) gives their actual
zero-inertia row. Thus these coordinates keep
the boundary topology of (E104.15)--(E104.16).

\emph{All comparison fibres, including zero inertia.}
Apply $N\to M$ to each actual cocycle value.
For the inert family the compatible $J$-adapted
Hall and Chen coordinates of E.100 delete exactly
the non-Chen negative coordinates of $y$.
This gives (E104.11) with both inverse maps.
For the split family the map is
$(D,r,b,[g])\mapsto(D,r,b,[\bar g])$.
Every fixed element has an actual fixed lift:
from $h_0$ lifting $\bar h\in M^J$, put
$k=h_0^{-1}J(h_0)\in N''$ and use $h_0k^{1/2}$.
Above a chosen lift $g_0$ of $\bar g$, the full
fibre map is
\[
 (N'')^J\backslash N''\longrightarrow\operatorname{Fib}([\bar g]),
                     \quad[k]\longmapsto[kg_0].
 \tag{E104.18}
\]
Removing the lifted fixed factor gives its inverse.
The integral kernel anti coordinates of E98.28--31
then give exactly (E104.11), with all nonlinear
central corrections retained.

With inertia value one, the cocycle kills the
entire closed normal inertia subgroup, because
the based action of that subgroup is trivial.
It therefore factors through the actual quotient
$\Gamma$, and every cocycle on $\Gamma$ inflates.
The new unramified free basis in either family is
transported to the original basis of E.41 by the
actual inverse free-basis maps above. Thus its
section quotient and every full-to-metabelian
fibre are exactly those of E100.22--26, with
the original $T$ in (E104.12). This comparison
does not replace $T$ by the prime-dependent
$F_\ell$ or $T_\ell$. E.102 evaluates its
first kernel factor by actual integral Smith maps.
Finally every ramified full class has nonzero
peripheral abelianization $r[d_D]$. The actual
quotient to $M$ preserves it, so every fibre
over zero inertia has already been included.

The retained program checks every prime through
$2000$ in the two sets (E104.1), the complete
two-by-two Frattini inverses, original-root
conjugation identities, all twenty-seven inert
descent labels, and the finite-precision forward
and inverse exponent maps in (E104.17).
It independently evaluates the Hall character
ranks through degree twenty and compares the
first twelve rows with the frozen actual
E.100 certificate. These checks support the
displayed generator and coordinate formulas;
the source isomorphisms, complete local
classification and all-degree inverse maps
above prove the statement for the full prime
sets (E104.1).
\end{proof}
